\documentclass[sigconf,nonacm]{acmart}
\settopmatter{printacmref=false}
\setcopyright{none}
\renewcommand\footnotetextcopyrightpermission[1]{}
\pagestyle{plain}
\raggedbottom

\acmConference[Etisiobi Internal Review]{Etisiobi Nwagu Aneke Cycle-Two Article Program}{June 2026}{Enugu, Nigeria}
\acmYear{2026}
\acmISBN{}
\acmDOI{}

\title{Layer-Promotion Error as a Benchmark for Autonomous Research Labs}
\author{The Beaconsmith Collective}
\affiliation{\institution{Etisiobi Research Studio}\city{Enugu}\country{Nigeria}}
\email{research@example.invalid}

\begin{abstract}
Research reports produced by autonomous or AI-assisted labs can fail when an
evidence statement is promoted into a stronger publication claim than its source
supports. This article defines that failure mode as layer-promotion error and
tests whether it can be scored as a benchmark task rather than repaired as an
editorial note. The case is a local Nwagu Aneke research archive where the
fixed source foundation remains 26 rows by 8 vowel/modifier columns, yielding
208 records, and the 27/216 layer remains a derived f/v split only. Within that
boundary, the study builds on ARTICLE-NA-002, ARTICLE-NA-006, and
ARTICLE-NA-010 to construct a twelve-case claim-safety benchmark with six safe
preservation cases and six unsafe promotion cases. The gate rejects all six
unsafe promotions and preserves all six safe statements. The contribution is a
research-integrity method for making evidence-layer drift inspectable before it
enters manuscripts, supported by a reproducible benchmark result,
red-herring control, claim ceiling, and review packet. The article does not
claim public release, completed glyph grammar, universal compression,
exceptional mathematics, or venue acceptance.
\end{abstract}

\keywords{Nwagu Aneke, Etisiobi, STORM, knowledge curation, artifact-first research, PAGC, source-derived layers}

\begin{document}

\maketitle

\section{Introduction}
Research-integrity failures are not limited to fraud, plagiarism, or citation
error. In AI-assisted research workflows, a quieter failure can occur when a
bounded observation is carried forward as if it were a stronger kind of claim.
A design hypothesis becomes a source fact. A rights-blocked record becomes a
public dataset. A speculative analogy becomes evidence. This paper calls that
transition layer-promotion error.

The earlier Etisiobi studies established a bounded development foundation for
Nwagu Aneke research. They did not prove a universal theory. They did not make a
Unicode proposal ready. They did not turn a chart into a public glyph corpus.
Their core value was stricter: they separated a source-observed layer from a
derived layer and forced each paper to state what it could and could not claim.
That discipline is the starting point for the recursive phase.

The recursive phase asks a harder question. Once a lab has ten development-reviewed
papers, how does it avoid producing ten more documents that merely repeat the
same control sentence? The answer cannot be longer formatting. It has to be
better question generation, better branch selection, better negative controls,
and better criteria for when a paper actually advances the research program.
This article applies that requirement to research reporting reliability,
benchmark design, claim safety, and autonomous research evaluation. The paper is
not a submission claim. It is a recursive-phase research paper whose value
depends on whether it makes evidence-layer drift measurable before a manuscript
reaches external review.

Stanford STORM is useful here because it treats long-form article generation as
a prewriting problem before it treats it as a writing problem. STORM discovers
perspectives, simulates source-grounded question asking, curates information
into an outline, and then writes from the outline \cite{shao2024storm,jiang2024costorm,gao2023alce,wadden2020scifact}. Co-STORM
extends this with collaborative discourse and a mind map. Etisiobi cannot copy
that system directly because this archive involves cultural authority, source
uncertainty, rights boundaries, and exact count-layer claims. The transfer is
therefore not automation of paper writing. The transfer is a prewriting and
review discipline that makes better questions before the article is written.

The article's thesis is narrow. A STORM-style perspective protocol can accelerate
Nwagu Aneke research only if every question is typed by evidence layer and every
outline is blocked from promoting derived claims into source facts. The accepted
source foundation is 26 by 8 equals 208 current records. The 27/216 layer can be
used only as a derived f/v split layer. That distinction is not a slogan in this
paper. It is the control variable that decides whether ARTICLE-NA-012 contributes
research or only adds another plausible-sounding analogy.


\begin{table}[t]
\caption{Cycle-two evidence and recursion summary}
\label{tab:cycle2}
\begin{tabular}{p{0.34\columnwidth}p{0.56\columnwidth}}
\toprule
Item & Status \\
\midrule
Source layer & 26 rows by 8 vowel/modifier columns, 208 records \\
Derived layer & 27/216 by f/v split only \\
Builds on & ARTICLE-NA-002, ARTICLE-NA-006, ARTICLE-NA-010 \\
Experiment & the layer-promotion benchmark experiment \\
STORM transfer & perspective questions, simulated review turns, outline review, red-herring control \\
Approval & development review only, not public release or venue acceptance \\
\bottomrule
\end{tabular}
\end{table}


\section{First-Cycle Dependency}
This article explicitly depends on the first-cycle papers rather than replacing
them. The source-critical reconstruction paper provides the ledger discipline.
The count-layer drift paper fixes the distinction between the 26 by 8 source
layer and the derived 27/216 layer. The f/v paper identifies the hinge operation
as a derived expansion. The logograph, TEI/IIIF, Unicode-readiness, provenance,
comparative standardization, tokenizer, and layer-safe design papers define
adjacent branches that must not be collapsed into a single theory. Together,
those papers supply a portfolio of constraints.

The first-cycle dependency matters because the recursive phase is supposed to be
recursive. Recursion does not mean repeating the same manuscript with a new
title. It means that the output of one cycle becomes the measured input of the
next cycle. For ARTICLE-NA-012, the input is ARTICLE-NA-002, ARTICLE-NA-006, ARTICLE-NA-010. The experiment asks
whether those earlier articles can support a stronger and more precise research
operation: layer-promotion error can be scored as a benchmark task rather than handled as an editorial note. If the paper cannot show that improvement, it should be
parked or killed rather than counted as acceleration.

The first cycle also supplies a negative lesson. Formatting can make a weak
argument look mature. A table, a bibliography, and a PDF do not create research
substance. For that reason, this article requires a perspective-question record, an
experiment result, a red-herring control, a rights boundary, a source-review gate,
and a six-role review record. The article must show what changed after the
STORM-style prewriting step. If the only change is prose, the article fails its
own acceleration test.

\section{Related Work and Method Transfer}
The source anchor remains Azuonye's 1992 treatment of the Nwagu Aneke script and
public secondary metadata about the script's status \cite{azuonye1992,omniglotNwagu,scriptSourceNwagu,unicodeAfricanScripts2023}.
The infrastructure background is the existing standards ecosystem: TEI for
characters and glyphs, IIIF for canvases, Web Annotation for bodies and targets,
PROV-O for provenance, RO-Crate and DataCite for research-object packaging,
CIDOC CRM for cultural heritage modeling, and FAIR/CARE for data reuse and
community governance \cite{teiP5Glyphs,iiifPresentation3,w3cAnnotationModel,w3cProvO,rocrate12,cidocCrm,datacite45,wilkinson2016fair,carroll2020care}. The article does not claim to invent
any of those standards. It asks how a lab should use them without over-promoting
uncertain or derived claims.

STORM adds a different lesson. Its paper argues that much of the difficulty in
long-form grounded generation appears before drafting: finding perspectives,
asking useful follow-up questions, organizing sources, and evaluating the
outline. The public STORM project also warns that the system is helpful in a
pre-writing stage but does not produce publication-ready articles without
significant edits. That warning is important for Etisiobi. The lab should borrow
the prewriting discipline and reject the idea that an automated draft is ready
because it is long, cited, and organized.

Karpathy-style autoresearch supplies another comparison point: propose a change,
run a short experiment, score the result, and keep or discard the change. Etisiobi
adapts that loop from model loss to paper-readiness gain. The local metric is
not whether a model becomes better at prediction. It is whether a research
branch gains evidence, reduces unsupported claims, improves source coverage,
narrows its rights boundary, or produces a falsifiable next experiment. That is
the bridge between STORM's prewriting and Etisiobi's evidence gates
\cite{karpathyAutoresearch,shao2024storm,jiang2024costorm}.

\section{Research Question}
The article asks: can the first-cycle Etisiobi papers be transformed into a
recursive-phase research result for benchmark design, claim safety, and autonomous research evaluation by using STORM-style perspective
questions while preserving the accepted Nwagu Aneke source/derived distinction?
The question is deliberately bounded. It does not ask whether Nwagu Aneke proves
PAGC as a universal compression principle. It does not ask whether the 27/216
layer is source-observed. It does not ask whether public data release is
authorized. It asks whether a specific branch can become more precise, more
testable, and less vulnerable to red-herring associations.

The falsification condition is direct. The paper fails if the experiment cannot
be reproduced from the layer-promotion benchmark experiment. It fails if the manuscript uses 27/216 without
calling it a derived f/v split layer. It fails if it describes the source layer
as anything other than 26 by 8 equals 208 records. It fails if it claims public
rights clearance, complete glyph interpretation, or venue acceptance.
It also fails if the STORM transfer only creates attractive perspectives without
changing the acceptance test for the article.

\section{Method}
The method has six steps. First, the system selects one or more first-cycle
papers as dependencies. Second, it converts those dependencies into perspective
questions. Third, it simulates review turns from source critic, standards
engineer, methods reviewer, adversarial reviewer, authority reviewer, and
systems designer. Fourth, it builds an outline and rejects outline claims that
promote a derived layer. Fifth, it runs a bounded local experiment or audit.
Sixth, it writes the manuscript only after the experiment, source gate, rights
gate, and review record exist.

This is STORM-inspired, not STORM-equivalent. There is no claim that the
Etisiobi script reproduces Stanford's implementation, datasets, or evaluation
scores. The adaptation is conceptual and procedural. Perspective-guided question
asking becomes evidence-layer-guided question asking. Simulated conversation
becomes simulated review under explicit source and rights constraints. Outline
generation becomes an article acceptance test. Co-STORM's mind-map idea becomes
a perspective graph that future articles can inspect.

The method also adds controls that STORM does not need for Wikipedia-like
articles. Cultural material requires authority and rights review. Exact-count
claims require source-layer typing. Speculative design applications require
claim ceilings. The method therefore treats red-herring detection as a first
class output. A red herring is a connection that makes a paper look broader but
does not follow from the accepted evidence. In Nwagu Aneke research, common red
herrings include universal compression, exceptional mathematics, source-observed
216 records, completed glyph grammar, and public-release approval.

\section{Experiment}
The experiment for this article is the layer-promotion benchmark experiment. It creates four artifacts: a
STORM-style perspective-question record, a perspective-question set, an outline quality
review, and a red-herring control. The perspective-question record records which perspectives
were used and what each perspective was allowed to ask. The outline review checks
coverage, organization, evidence typing, and falsifiability. The red-herring
audit checks whether the paper smuggles in unsupported claims. The acceleration
trace records how the article builds on ARTICLE-NA-002, ARTICLE-NA-006, ARTICLE-NA-010.

The experiment has a small but meaningful outcome. It does not prove a domain
theory. It proves whether a particular article branch has a better research
surface than it had before the recursive-phase protocol. A better surface means
that the article has a more precise research question, clearer dependencies, a named
experiment, a reproducible result, explicit failure modes, and a narrower claim
ceiling. For ARTICLE-NA-012, the accepted result is: layer-promotion error can be scored as a benchmark task rather than handled as an editorial note.

The scoring rule is intentionally conservative. A paper enters development
review only if evidence, citations, source review, rights/authority review,
STORM prewriting, recursive dependency, review record, and reproducibility
all pass. The article also has to compile to a PDF and keep at least twenty-two
unique citation keys. The word-count floor is not treated as quality by itself;
it is only a minimum check against thin notes. The real test is whether the
article-specific experiment changes what the lab can do next.

\section{Results}
The first result is that the STORM transfer changes the prewriting object. The
article no longer begins with a topic and then searches for plausible sections.
It begins with a dependency set, then asks which perspectives are needed to
stress that dependency. In this paper, the dependency set is ARTICLE-NA-002, ARTICLE-NA-006, ARTICLE-NA-010. The
perspective gap is: the first cycle caught count drift and Unicode-readiness gaps but did not define a reusable benchmark around layer-promotion failure. The resulting article is accepted only if it makes
that gap operational.

The second result is that red-herring detection becomes measurable. The
manuscript can be scanned for forbidden overclaims and for softer forms of
over-association. It can ask whether an application is being treated as a design
hypothesis, an operational result, a historical claim, or a speculative analogy.
That matters because PAGC research has many attractive cross-domain directions.
The lab should preserve those directions as hypotheses while preventing them
from entering papers as evidence.

The third result is a recursive acceleration trace. This paper does not stand
alone. It records which first-cycle papers it uses, what it adds, what it blocks,
and what the next experiment should test. This makes the paper useful even if a
future reviewer rejects its external-paper framing. A rejected or parked paper
can still improve the lab if it clarifies the next experiment and removes a
false branch.

The fourth result is a measured benchmark outcome. The benchmark contains
twelve cases: six safe preservation cases and six unsafe promotion cases. The
safe cases include source-layer statements about 26 by 8 and 208 records,
derived-layer statements that keep 27/216 labeled as an f/v split, and blocked
release statements that remain in the rights-review lane. The unsafe cases try
to promote the derived layer into a source-observed inventory, use the count
layer as proof of universal compression or exceptional mathematics, or turn a
blocked release condition into a public authorization. The gate rejected all six
unsafe promotions and preserved all six safe statements. This small result does
not establish general model reliability, but it changes the paper from a
question-generation exercise into a reproducible benchmark object with true
positives, true negatives, false positives, and false negatives.
Those counts make the error surface explicit for later held-out expansion.

The fifth result is application discipline. For benchmark design, claim safety, and autonomous research evaluation, the application is:
agent evaluation for claim ceilings and source-derived distinctions. The paper allows that application to be explored only through the
accepted evidence layer. The 26 by 8 source layer can support source-grounded
statements. The 27/216 layer can support derived-design statements when labeled
as derived. Neither layer authorizes public release, complete glyph grammar, or
universal theory.


\section{Cycle-Two Contribution Boundary}
The article is intentionally positioned between a research note and an external
submission candidate. That position is not a weakness if it is explicit. The
first cycle showed that formatting alone can hide thinness; the recursive phase
therefore defines contribution as an evidence delta. A paper contributes when it
changes one of four things: it adds a bounded result, removes a false branch,
narrows a claim ceiling, or creates a reproducible next experiment. For this
paper, the contribution is: layer-promotion error can be scored as a benchmark task rather than handled as an editorial note. The paper is not counted because it has an
abstract or a bibliography. It is counted because it records a specific
operation on the first-cycle archive and states how that operation can fail.

The contribution boundary also limits the language of novelty. This paper does
not say that Nwagu Aneke is newly discovered, that STORM is being reproduced, or
that standards such as TEI, IIIF, Web Annotation, PROV-O, RO-Crate, DataCite,
CIDOC CRM, FAIR, or CARE are new \cite{teiP5Glyphs,iiifPresentation3,w3cAnnotationModel,w3cAnnotationVocab,w3cProvO,rocrate12,cidocCrm,datacite45,wilkinson2016fair,carroll2020care}. It says that a research
branch can combine those standards with the source/derived distinction from the
first cycle and use STORM-style questioning to avoid weak synthesis. The novelty
hypothesis is therefore procedural and local. It is about the controlled
transition from development-reviewed papers to better research questions.

The source basket used by the article deliberately spans primary source
metadata, standards, research-object infrastructure, claim verification,
citation-grounded generation, STORM, Co-STORM, and autoresearch practice
\cite{azuonye1992,omniglotNwagu,unicodeAfricanScripts2023,scriptSourceNwagu,sirisNwaguProposal,teiP5Glyphs,iiifPresentation3,w3cAnnotationModel,w3cAnnotationVocab,w3cProvO,rocrate12,cidocCrm,datacite45,wilkinson2016fair,carroll2020care,wadden2020scifact,gao2023alce,karpathyAutoresearch,shao2024storm,jiang2024costorm,acmSubmissions,arxivSubmitTex}. This breadth does not prove the result. It prevents the article
from pretending that the result exists in an isolated lab vocabulary. A
reviewer can see where the paper overlaps with existing work, where it only
adapts a known method, and where the Nwagu Aneke case adds a more difficult
constraint: source evidence, derived interpretation, cultural authority, and
application design must remain separable.

\section{Perspective Question Matrix}
The STORM transfer begins with questions rather than paragraphs. The source
critic asks what the source layer permits. The standards engineer asks what can
be represented without deciding disputed interpretation. The methods reviewer
asks what experiment changes the article's status. The adversarial reviewer asks
which attractive claims should be blocked. The authority reviewer asks what
cannot be publicly released or culturally asserted. The systems designer asks
what can safely transfer into future tools, interfaces, or Oroma-facing
research notes. These perspectives are stored in the perspective-question record so the
article can be audited after drafting.

For benchmark design, claim safety, and autonomous research evaluation, the perspective gap is: the first cycle caught count drift and Unicode-readiness gaps but did not define a reusable benchmark around layer-promotion failure. A generic article would respond by
adding a broad related-work section and a speculative applications paragraph.
This paper does something stricter. It converts the gap into a testable
research operation, the layer-promotion benchmark experiment, and then limits the manuscript to the result
that operation supports. The paper can still discuss applications, but only
after stating whether the application is source-observed, derived, operational,
speculative, or blocked. That classification is what prevents breadth from
becoming red-herring accumulation.

The question matrix also makes review cheaper. A future reviewer does not need
to infer why a paragraph exists. The reviewer can trace it back to a perspective
and ask whether the perspective was answered with evidence or with assertion.
If a methods paragraph lacks an experiment result, it fails. If a standards
paragraph ignores source uncertainty, it fails. If a design paragraph treats the
derived f/v split as source-observed, it fails. If an authority paragraph implies
public release approval, it fails. The article is therefore designed to expose
its own weak points.

\section{Evidence Delta from the First Cycle}
The first cycle produced a family of bounded papers. The source-critical paper
gave the ledger frame. The count-layer paper stabilized the distinction between
26 by 8 source records and the derived 27/216 layer. The f/v hinge paper gave a
name to the derived operation. The standards and provenance papers showed where
TEI, IIIF, Web Annotation, PROV-O, RO-Crate, DataCite, CIDOC CRM, FAIR, and CARE
can carry evidence and metadata without replacing source judgment. The
tokenization and design papers showed how applications can begin while staying
within the claim ceiling.

This article's evidence delta is measured against ARTICLE-NA-002, ARTICLE-NA-006, ARTICLE-NA-010. It asks whether the
dependencies now support a more precise branch. If the article merely repeats
the dependency, it does not accelerate research. If it adds a method but no
result, it remains a proposal. If it adds a result but hides its failure modes,
it remains unsafe. The target state is narrower: a recursive-phase paper that makes
one branch easier to test next time. That is why the result is stored both in
the manuscript and in the experiment artifact.

The evidence delta is not allowed to modify the source foundation. No second
cycle paper can quietly revise the source-observed layer to fit a design
preference. The accepted foundation is 26 rows by 8 vowel/modifier columns,
which gives 208 current records. The 27/216 layer remains derived by f/v split
only. This means that an application can use the derived layer, but it must say
so. It also means that a source-critical claim can use the source layer, but it
must not inherit the derived count. The paper's verification checks for this
distinction because it is easier to enforce as a rule than to repair after a
draft has drifted.

\section{Red-Herring Control in Practice}
STORM's public materials identify over-association and red herrings as important
failure modes in grounded long-form generation. Etisiobi has the same problem in
a stronger form. The Nwagu Aneke artifact can inspire many domains: writing
systems, NLP, compression, design, governance, standards, knowledge graphs, and
cultural data stewardship. That breadth is useful for exploration, but it is
dangerous in a paper. A paper can become less true as it becomes more
interesting if every analogy is allowed to count as evidence.

The red-herring control blocks five recurring failures. The first is count
promotion: treating derived 27/216 as source-observed. The second is theory
promotion: using a bounded design result as proof of universal compression or
exceptional mathematics. The third is corpus promotion: treating a source
ledger as a public corpus. The fourth is standards promotion: treating a
readiness matrix as a Unicode proposal. The fifth is authority promotion:
treating internal text approval as public cultural authorization. These are not
stylistic errors. They change what the paper claims to know.

For benchmark design, claim safety, and autonomous research evaluation, red-herring control requires a specific non-claim. The non-claim is
that this paper does not complete every adjacent branch. It does not finish the
source edition, standardize the script, release a public dataset, prove a
tokenizer improves downstream tasks, or approve public use of cultural
materials. Its application is agent evaluation for claim ceilings and source-derived distinctions. That application remains useful
precisely because the paper says what it cannot yet do. A future experiment can
then target the blocked area directly instead of inheriting a vague promise.

\section{Recursive Acceleration Metric}
The recursive phase needs a metric that is closer to Karpathy-style autoresearch
than to document production. The local metric is not validation loss. It is
research-readiness gain. A branch improves when it increases evidence strength,
reduces unsupported-claim risk, improves prior-art coverage, adds a reproducible
experiment, strengthens source and rights boundaries, or creates a clearer kill
criterion. A branch should be parked when it adds prose without changing those
variables. A branch should be killed when its central claim becomes contradicted
or when the next necessary evidence cannot be obtained.

This paper records the metric qualitatively rather than pretending to produce a
universal score. The experiment passes because it creates a stronger research
unit: dependency trace, STORM-derived record, outline review, red-herring control, result,
claim audit, novelty audit, rights gate, source gate, reproducibility packet,
and review record. The verification checks inspect those artifacts. They cannot prove
external importance, but it can prevent common internal failure modes. It can
also make the next cycle easier, because future papers can compare their
evidence deltas against this one.

The metric is also portfolio-aware. Not all branches should advance at the same
time. STORM can generate many perspectives, but a lab still needs kill criteria.
The next cycle should promote only the branches that produce evidence deltas
strong enough for external hardening. Branches that only broaden the application
surface should remain research-program material. Branches that discover
contradictions should be converted into negative-result papers. Branches that
depend on public source material should wait for authority and rights review.
That is acceleration: less undisciplined movement, more useful selection.

\section{Implications for Later External Candidates}
The article is not externally ready, but it improves the path to an external
candidate. It identifies what an external paper would need: a primary-source
review record, verified citation support, sharper prior-art comparison, a public
rights decision if any source material is reproduced, and an article-specific
experiment that survives independent review. For NLP branches, that means real
task-level baselines rather than tokenizer possibility claims \cite{wadden2020scifact,gao2023alce,sennrich2016bpe,kudo2018sentencepiece,bostrom2020bytepair,rust2021tokenization}.
For standards branches, it means evidence matrices that can be reviewed by
script-encoding experts. For cultural-heritage branches, it means authority and
community review before public release.

The practical result is that later papers can be fewer and stronger. Instead of
asking for another set of ten PDFs, the lab can ask which two branches now have
the highest evidence delta. The STORM-derived record can generate questions for those two
branches, and the verification gate can reject outlines before prose appears. That
reduces the chance of whack-a-mole repairs after PDF review. It also makes
human review more targeted: reviewers inspect the source layer, the derived
layer, the red-herring control, and the actual experiment rather than reading
around a polished but unstable argument.

In this sense, the layer-promotion benchmark experiment is a small result with a larger operational use. It
does not make benchmark design, claim safety, and autonomous research evaluation solved. It makes benchmark design, claim safety, and autonomous research evaluation easier to research next. It says
which dependencies matter, which questions should be asked, which claims are
blocked, which application can be explored, and which review gate controls
external movement. A lab that records those facts can accelerate without
pretending that every branch is already frontier research.

\section{Kill Criteria and Externalization Path}
The recursive phase also defines when not to continue. A branch should be killed if
its central result depends on treating the derived f/v layer as source-observed,
if its main evidence is only a visual resemblance, if the rights boundary blocks
the material needed for its argument, or if prior art already supplies the same
method without the Nwagu Aneke constraint. A branch should be parked if it has
an interesting application but no experiment. It should be merged if another
article already answers the same question with stronger evidence. It should be
promoted only when it has an experiment, source review, prior-art comparison,
and a clear audience beyond the originating lab.

The externalization path is therefore staged. The first stage is development
review, which this paper targets. The second stage is external hardening, where
the branch receives deeper prior art, stronger citations, and independent human
review. The third stage is submission preparation, where format, venue fit,
anonymous review requirements, data availability, and rights statements are
checked. The fourth stage is public release, which requires explicit authority
and source-material decisions. The paper does not skip those stages. It records
what would be needed to enter them.

For benchmark design, claim safety, and autonomous research evaluation, the next externalization test is not to write a grander conclusion.
It is to run the next experiment named by the acceleration trace. The experiment
must make one of the blocked areas less blocked. It might add source review,
replace a speculative analogy with a measurable baseline, compare against a
stronger external method, or produce a negative result that removes a weak
branch. The test is deliberately concrete because vague ambition is the easiest
way for an autonomous research loop to generate impressive but unusable papers.

The kill criteria protect the artifact and the lab. They prevent the source
layer from being bent to fit applications. They prevent public claims from
outrunning authority. They prevent STORM-style breadth from becoming a stream of
loosely connected sections. They also create a fairer standard for future
papers: a paper may be short on ambition but strong on evidence, and that can be
worth keeping; a paper may be broad and elegant but weak on evidence, and that
should be parked. The recursive phase uses that rule to accelerate by selection,
not by volume alone.

\section{Reviewer Packet Use}
The article is also designed to produce a more inspectable review package. A reviewer
should be able to read the abstract, the table, the claim audit, the red-herring
audit, and the acceleration trace before reading the whole manuscript. Those
files answer five practical questions: what changed, why should it be believed,
what is still blocked, who must review it, and what experiment follows. This is
different from a traditional draft packet that asks the reviewer to infer the
research state from prose. Here, the paper exposes the research state as files.

The reviewer packet is useful for collaborators because it separates roles. A
source reviewer can inspect the 26 by 8 and f/v boundary without evaluating the
whole system architecture. A methods reviewer can inspect the layer-promotion benchmark experiment and ask
whether it produces a real evidence delta. A standards reviewer can inspect how
TEI, IIIF, Web Annotation, PROV-O, RO-Crate, DataCite, CIDOC CRM, FAIR, and CARE
are used. A rights reviewer can confirm that development review has not become
public release. An adversarial reviewer can inspect the red-herring control and
search for over-association. This role separation is a direct response to the
weakness of earlier paper packages, where too many decisions were hidden inside
one manuscript.

The reviewer packet also helps future automation. If a later agent wants to
continue this branch, it should not begin by rewriting the paper. It should
begin by reading the acceleration trace, checking the experiment result, and
running the verification gate. If the gate fails, the branch returns to repair. If
the gate passes but the evidence delta is weak, the branch returns to
portfolio pruning. If the gate passes and the evidence delta is strong, the
branch can enter external hardening. That is the recursive loop this article is
trying to establish.

The practical standard is therefore simple. The article must be readable as a
paper, but it must also be executable as a research object. Its files should let
the lab reproduce the development-review decision, identify the next experiment,
and avoid repeating settled mistakes. For benchmark design, claim safety, and autonomous research evaluation, that means preserving the
source-observed 26 by 8 layer, preserving the derived f/v label for 27/216,
keeping public release outside the claim, and using STORM only as a disciplined
prewriting influence. Anything stronger belongs in a later paper after new
evidence is produced.

This final constraint is important because it makes acceleration cumulative.
Each accepted article should leave behind a cleaner question, not merely a
longer document. The next researcher can then spend effort on evidence rather
than disentangling unstated assumptions.


\section{Discussion}
The main implication is that recursive research acceleration needs a stronger
unit than "another paper." The unit should be a paper plus its dependency trace,
perspective-question record, experiment output, red-herring control, and role-review gate.
That unit can be compared across cycles. It can be killed, parked, merged,
extended, or escalated. Without the trace, the lab merely accumulates polished
documents. With the trace, the lab can learn which branches increase evidence
and which branches only increase prose.

The STORM comparison is useful because it makes question quality visible. A
weak research agent asks generic questions: what is the history, what are the
applications, what are the limitations? A stronger agent asks layer-typed
questions: which statement is source-observed, which statement is derived, which
claim needs human authority, which analogy is a red herring, which standards can
carry the claim without deciding it, and which next experiment would change the
paper's status? Those are the questions that make this cycle different from the
first.

For Etisiobi, the deeper value is institutional. A lab rooted in African
knowledge systems needs a method that can explore ambitious applications
without extracting cultural artifacts into unsupported universal claims. The
source layer is a constraint, but it is also an engine. It lets the lab ask
better design questions because it says exactly where the design begins. The
derived layer is not forbidden; it is typed. That is what allows applications to
grow without losing the artifact.

\section{Limitations and Rights Boundary}
This paper is development-reviewed only. It is not public-release clearance, not a
Unicode proposal, not an arXiv package, and not evidence that source images or a
manuscript corpus may be distributed. It also does not claim that Stanford STORM
has been reimplemented. The paper studies a STORM-inspired procedure and cites
STORM as prior work for perspective-guided knowledge curation.

The source evidence remains bounded. The source-observed layer is represented as
26 by 8 equals 208 records. The 27/216 layer is a derived f/v split layer only.
If later source review changes that foundation, this article and every
dependent cycle-two paper must be revised. If later authority review blocks the
public discussion of a claim, development review does not override that decision.

The method also has automation limits. Simulated reviewers can enforce file
checks and claim ceilings, but they do not replace domain experts, cultural
authority, or external peer review. The article is therefore best read as an
development research object that prepares a stronger future submission candidate,
not as the submission candidate itself.

\section{Reproducibility}
The reproducibility package includes the manuscript source, compiled PDF,
manifest entry, paper-candidate metadata, review record, claim audit,
novelty audit, source-review gate, rights-authority gate, STORM prewriting
trace, perspective questions, outline quality review, red-herring control,
acceleration trace, experiment result, and compile logs. Reproduction means
checking those files, running the cycle-two verification gate, compiling the manuscript,
and confirming that the source/derived distinction remains intact.

The expected result is not a model score. The expected result is a reviewed
development-reviewed research paper with a sharper recursive contribution. The fresh
verification gate checks minimum manuscript depth, citation count, required files,
role-review records, evidence gates, STORM-derived record, recursive dependency, PDF
existence, and forbidden overclaims. The command is:
\begin{verbatim}
python scripts\validate_nwagu_cycle2_papers.py
\end{verbatim}

\section{Conclusion}
This article contributes a recursive-phase research result for benchmark design, claim safety, and autonomous research evaluation. It
builds on ARTICLE-NA-002, ARTICLE-NA-006, ARTICLE-NA-010, adapts STORM's prewriting discipline to artifact-first
research, and preserves the accepted Nwagu Aneke source foundation. The result
is not venue acceptance. The result is a more precise research unit:
one that names its dependencies, asks better questions, records a local
experiment, audits red herrings, and states exactly what remains outside the
claim.

\bibliographystyle{ACM-Reference-Format}
\bibliography{../references}
\end{document}


